\documentclass[10pt,a4paper]{amsart}
\usepackage[margin=19mm]{geometry}
\usepackage{amsmath,amssymb,mathtools}
\usepackage[scr=rsfs,cal=euler]{mathalfa}
\usepackage{xfrac}
\usepackage[unicode,hidelinks,bookmarks=false]{hyperref}
\newcommand{\ie}{i.e.\ }
\newcommand{\cf}{cf.\ }
\newcommand{\resp}{respectively\ }
\newcommand{\Subsection}{\subsection}
\providecommand{\nobreakdash}{\nobreak\hbox{-}}
\setlength{\emergencystretch}{2em}
\multlinegap0pt
\newcommand{\Ps}{\mathbf{P}}
\newcommand{\As}{\mathbf{A}}
\newcommand{\C}{\mathbf{C}}
\newcommand{\Q}{\mathbf{Q}}
\newcommand{\Z}{\mathbf{Z}}
\newcommand{\F}{\mathbf{F}}
\newcommand{\R}{\mathbf{R}}
\newcommand{\G}{\mathbf{G}}
\newcommand{\cL}{\mathcal{L}}
\newcommand{\cO}{\mathcal{O}}
\newcommand{\cD}{\mathcal{D}}
\newcommand{\cI}{\mathcal{I}}
\newtheorem{lemma}{Lemma}[section]
\newtheorem{proposition}[lemma]{Proposition}
\newtheorem{theorem}[lemma]{Theorem}
\newtheorem{corollary}[lemma]{Corollary}
\newtheorem{conjecture}[lemma]{Conjecture}
\newtheorem{notation}[lemma]{Notation}
\newtheorem{definition}[lemma]{Definition}
\newtheorem{construction}[lemma]{Construction}
\newtheorem{assumption}[lemma]{Assumption}
\newtheorem{conv}[lemma]{Convention}
\newtheorem{remark}[lemma]{Remark}
\newtheorem{example}[lemma]{Example}
\DeclareMathOperator{\degsp}{degsp}
\DeclareMathOperator{\CH}{CH}
\DeclareMathOperator{\NL}{NL}
\DeclareMathOperator{\red}{red}
\DeclareMathOperator{\prim}{prim}
\DeclareMathOperator{\codim}{codim}
\DeclareMathOperator{\rank}{rank}
\DeclareMathOperator{\rk}{rk}
\DeclareMathOperator{\Aut}{Aut}
\DeclareMathOperator{\Cl}{Cl}
\DeclareMathOperator{\sat}{sat}
\DeclareMathOperator{\Div}{Div}
\DeclareMathOperator{\Pic}{Pic}
\DeclareMathOperator{\length}{length}
\DeclareMathOperator{\Gr}{Gr}
\DeclareMathOperator{\Hom}{Hom}
\DeclareMathOperator{\NS}{NS}
\DeclareMathOperator{\cha}{char}
\DeclareMathOperator{\MW}{MW}
\DeclareMathOperator{\ima}{im}
\DeclareMathOperator{\sing}{sing}
\DeclareMathOperator{\spe}{sp}
\begin{document}
\title[Nodal deformations of hypersurfaces with an ordinary m-fold ]{Nodal deformations of hypersurfaces with an ordinary $m$-fold point}
\author{Remke Kloosterman}
\date{}
\maketitle
\noindent\emph{Source and licence.} Nodal deformations of hypersurfaces with an ordinary $m$-fold point. Version arXiv:2609.18313v1 (CC BY 4.0). This material is distributed under the Creative Commons Attribution 4.0 International licence, http://creativecommons.org/licenses/by/4.0/, as recorded in the arXiv OAI licence metadata for this exact version and shown on the abstract page https://arxiv.org/abs/2609.18313v1. Full rights notice: licenses/arxiv-2609.18313-CC-BY-4.0.txt.\par\smallskip
\noindent\textit{Editorial scope.} Counted unit: the original \texttt{theorem} environment \texttt{thmLowerSur} at source lines 323--328 together with its complete proof at lines 330--341; the delivered document keeps source lines 293--342 so that every referenced label and every citation resolves inside it. Native editable source is the paper's own concatenated TeX; the AMS wrapper is editorial formatting only. No publisher class, upstream script or unrelated artwork is redistributed.
\section{Lower bound for $\delta_S$}\label{secLower}
Let $S\subset \Ps^3$ be a surface of degree $d$ with a unique singular point, which is an ordinary $m$-fold point.

In \cite{CiGa26} it is shown that $\delta_S\geq \binom{m-1}{2}=\frac{m^2-m}{2}$. In the previous section we presented an upper bound which is cubic in $m$. 
In this section we show that  $\delta_S\geq \frac{m^3}{27}$.
%If one could extend the result \cite[Theorem 3.3]{CiGa26} to higher dimension, then this would yield a lower $\delta_X\geq  \frac{m^n}{n^n}$, for a hypersurface $X\subset \Ps^n$ of degree $d\geq m$ containing an ordinary $m$-fold point and no further singularities.

We start by a recalling a result from commutative algebra:

\begin{proposition} Let $\Sigma\subset \Ps^n$ be a zero-dimensional scheme, let $I=I(\Sigma)\subset \C[x_0,\dots,x_n]$ be the ideal of $\Sigma$ and let 
\[ 0 \to \oplus_{j=1}^{i_t} S(-a_{tj}) \to \dots \to \oplus_{j=0}^{i_1}S(-a_{1j}) \to S \to S/I \to 0\]
be a free graded resolution of $S/I$, where $S=\C[x_0,\dots, x_n]$ and the $a_{ij}$ are positive integers.
Then $H^1(\Ps^n,\cO(d)\otimes \cI_{\Sigma})=0$ for all $d\geq \max(a_{ij})-n$.
\end{proposition}

\begin{proof} Since $H^i(\Ps^n,\cO(d))=0$ for all $i,d>0$ and $H^0(\Ps^n,\cO\times \cI_{\Sigma})=I_d$ we obtain that
\[ 0 \to H^0(\cO(d)\otimes \cI_{\Sigma}) \to H^0(\cO(d))) \to H^0(\cO_{\Sigma}) \to H^1(\cO(d)\otimes \cI_{\Sigma}) \to 0 \] is exact.  
Hence  $H^1(\Ps^n,\cO(d)\otimes \cI_{\Sigma})=0$ if and only if $h_I(d)=\length \Sigma$. i.e., the Hilbert function and Hilbert polynomial of $I$ agree at $d$.

One has that $\dim S_k=\frac{(k+n)(k+n-1)\dots (k+1)}{n!}$ for $j\geq -n$. Hence the Hilbert polynomial and Hilbert function of $\Ps^n$ agree for $k\geq -n$. From this it follows directly that the Hilbert polynomial and Hilbert function of $S/I$ agree for all $d$ at least $\max(a_{ij})-n$.
\end{proof}

\begin{proposition}\label{prpHilbBound} Let $\Sigma\subset \Ps^n$ be a zero-dimensional scheme, which is a scheme-theoretic complete  intersection of multidegree $(d_1,\dots,d_n)$. 
Then $H^1(\Ps^n,\cO(d)\otimes \cI_{\Sigma})=0$ for all $d\geq (\sum_{i=1}^n d_i)-n$.
\end{proposition}
\begin{proof}
Let $F_1,\dots,F_n$ be  homogeneous generators of $I$. Then the Koszul complex on $F_1,\dots,F_n$ is a resolution of $S/I$. In particular, for this resolution one has that $a_{ij}\leq \sum_{k=1}^n d_k$ for all $i$ and $j$. The result now follows from the previous result.
\end{proof}



\begin{theorem}\label{thmLowerSur} Let $n,d\geq 3$ be integers.  Moreover, suppose $(n,d)\neq (3,3)$. 

There exists a nodal hypersurface $X\subset\Ps^{n}$ of degree $d$ with precisely $(\lfloor \frac{d}{n} \rfloor+1)^{n}$ nodes such that $H^1(I_{\Sigma} (d))=0$, where $\Sigma=X_{\sing}$. If, moreover, $n$ does not divide $d$, then  $H^1(I_{\Sigma} (d-1))=0$.

If $n$  divides $d$ then there exists a nodal hypersurface $X\subset\Ps^{n}$ of degree $d$ with precisely $\frac{d}{n} (\frac{d}{n}+1)^{n-1}$ nodes such that $H^1(I_{\Sigma} (d-1))=H^1(I_{\Sigma} (d))=0$. 
\end{theorem}

\begin{proof}
Fix $n$ general forms $F_1,\dots,F_n$ in $x_0,\dots,x_n$ of degree $\lfloor \frac{d}{n} \rfloor+1$. Let $I=(F_1,\dots,F_{n})$. Then $I$ is a complete intersection ideal and its vanishing locus in $\Ps^n$ is finite.
Consider the ideal  $I^2=(F_iF_j)_{1\leq i\leq j\leq n}$. Pick a general 
\[F=\sum_{1\leq i\leq j \leq n+1} G_{ij}F_iF_j \in (I^2)_d\] and let $X=V(F)$.
The generators of $I^2$ have  degree  at most $2\lfloor d/n \rfloor+2$. This is at most  $d$ for our choices of $n$ and $d$.
Hence the base locus of $(I^2)_d$ is precisely $V(I)$ and by Bertini's theorem $X$ is smooth outside $V(I)$, hence $\Sigma=X_{\sing}=V(I)$. It remains to check that each singular point is a  node. For each point in $\Sigma$, the dehomogenization of the  polynomials $F_1,\dots F_{n}$ yields a  system of local coordinates. Recall that matrix  $\frac{ \partial^2 F}{\partial F_i\partial F_j}$ equals the matrix with entries $G_{ij}$ for $i<j$, $G_{ji}$ for $i>j$ and $2G_{ii}$ for $i=j$.
In order to have an node at every point in $\Sigma$ we need that the determinant of this matrix is invertible modulo $I$, which holds true for a generic choice of the $G_{ij}$.

From Proposition~\ref{prpHilbBound} it follows that $H^1(I_{\Sigma}(k))=0$ for $k\geq n \lfloor \frac{d}{n} \rfloor $. Since $n \lfloor \frac{d}{n} \rfloor\leq d$. This proves the claim on $H^1(I_\Sigma(d))$.
 
 If $H^1(I_{\Sigma}(d-1))\neq 0$ then equality holds in the previous formula, in particular, $n \lfloor \frac{d}{n}\rfloor =d$, which is in turn equivalent to  $n\mid d$. Hence $H^1(I_{\Sigma}(d-1))=0$ if $n\nmid d$. In the case $d\mid n$  we replace $F_{n}$ by a form of one degree less and obtain also  $H^1(I_{\Sigma}(d-1))=0$. 
\end{proof}

\begin{thebibliography}{99}
\bibitem[CiGa26]{CiGa26} C.~Ciliberto and C.~Galati. \newblock On nodal deformations of singular surfaces in $\mathbb{P}^3$. \newblock Preprint, available at \texttt{https://arxiv.org/abs/2602.09177}, 2026.
\end{thebibliography}
\end{document}
